%-------------------------------------------------------------
%Dans cet exercice, les deux parties sont indépendantes.

%Une entreprise développe une intelligence artificielle (IA) capable de reconnaître des objets sur des images.

\subsubsection*{Partie A}

%On entraîne l'IA à partir d'une base de données de \np{50000} images réparties en 4 catégories : \og Objets du quotidien \fg, \og Animaux \fg, \og Véhicules \fg, \og Autres \fg.

%L'intelligence artificielle est testée pour mesurer sa précision et son efficacité. Les images sont réparties comme suit :

%\vspace{6pt}
%\begin{center}
%\begin{tabular}{|>{\columncolor{gris}}c|l|c|}\hline
%\rowcolor{gris}
%& \textbf{Type d'image} & \textbf{Nombre d'images} \\\hline
%2 & Objets du quotidien & 28\,000 \\\hline
%3 & Animaux & \np{12000} \\\hline
%4 & Véhicules & \np{8000} \\\hline
%5 & Autres & ? \\\hline
%\end{tabular}
%\end{center}

%\medskip

\begin{enumerate}
	\item $\np{50000}-(\np{28 000}+\np{12000}+\np{8000})=\np{2000}$ \\
	Il y a $\np{2000}$ images dans la catégorie \og Autres \fg{}.
	\item $\dfrac{90}{100}\times \np{28000}=\np{25200}$

	Il y a $\np{25200}$ images reconnues correctement dans cette catégorie.
	\item $\dfrac{\np{5600}}{\np{8000}}\times 100=70$.

	Donc 70\% des images dans la catégorie \og Véhicules \fg{} sont reconnues correctement. 
	\item $\dfrac{\np{28000}}{\np{50000}}=0,56$.
	
	La probabilité que l'image tirée soit l'image d'un \og Objet du quotidien \fg{} est 0,56.
\end{enumerate}

\bigskip

\subsubsection*{Partie B}

%L'intelligence artificielle, très utilisée dans le monde entier, nécessite une quantité importante d'électricité. L'énergie consommée peut s'exprimer en wattheures (Wh).
%
%En 2024, sa consommation annuelle est estimée à 82\,000 Gigawattheures (GWh).
%
%En comparaison, un collège consomme en moyenne 200\,000 kilowattheures (kWh) par an.
%
%\medskip
%
%\begin{center}
%\begin{tabularx}{0.35\linewidth}{|X|}\hline
%\textbf{Rappels :}\\
%-- 1\,kWh $= 10^3$\,Wh\\
%-- 1\,GWh $= 10^9$\,Wh\\ \hline
%\end{tabularx}
%\end{center}
%
%\medskip

\begin{enumerate}[resume]
	\item IA : $8,2\times 10^{13}$ Wh\\
	Collège : $2\times 10^8$ Wh
	\item $\dfrac{8,2\times 10^{13}}{2\times 10^8}=4,1\times 10 ^5$
	
	On pourrait alimenter \np{410000} collèges pendant un an avec la consommation électrique de l'intelligence artificielle.
	\item $\dfrac{\np{410000}}{\np{7100}}=\dfrac{4100}{71}\approx 57,74$

	On pourrait alimenter tous les collèges français avec la consommation électrique annuelle de cette intelligence artificielle pendant près de 58 ans.
\end{enumerate}


%-------------------------------------------------------------
